\documentclass[10pt,a4paper]{article}

\usepackage[margin=20mm]{geometry}
\usepackage{amsmath,amssymb,amsthm}

\theoremstyle{definition}
\newtheorem*{problem*}{Problem}

\setlength{\parindent}{0pt}
\setlength{\parskip}{5pt}
\newcommand{\Erho}{\mathbb E_{\rho}}
\newcommand{\Eq}{\mathbb E_{q}}
\begin{document}

\begin{center}
  {\Large\bfseries A Conserved Concentration Coupled to Surface Diffusion}
\end{center}

\section{Definitions}

Fix $p>6$ and $\eta\in(0,\tfrac12)$. Let
$\Sigma\subset\mathbb R^3$ be a closed, connected, oriented $C^\infty$
surface with unit normal $\nu_\Sigma$.

Let $\mathcal U\subset W_p^{4-4/p}(\Sigma)$ be an open set such that, for
every $\rho\in\mathcal U$,
\begin{equation}
  X_\rho:\Sigma\longrightarrow\mathbb R^3,
  \qquad
  X_\rho(x)=x+\rho(x)\nu_\Sigma(x),
\end{equation}
is an embedded normal graph and its induced unit normal $\nu_\rho$, chosen
compatibly with $\nu_\Sigma$, satisfies
\begin{equation}
  \nu_\Sigma\cdot\nu_\rho>0
  \qquad\text{on }\Sigma.
\end{equation}
The boundary $\partial\mathcal U$ is understood with respect to the topology
of $W_p^{4-4/p}(\Sigma)$.

Let
\begin{equation}
  \rho_0\in\mathcal U,
  \qquad
  \Gamma_0=X_{\rho_0}(\Sigma).
\end{equation}
Let $q_0$ be a scalar field on $\Gamma_0$ such that
\begin{equation}
  \widetilde q_0:=q_0\circ X_{\rho_0}
  \in W_p^{2-2/p}(\Sigma)
\end{equation}
and
\begin{equation}
  \eta\leq q_0\leq1-\eta
  \qquad\text{on }\Gamma_0.
\end{equation}

For an oriented evolving surface $\Gamma(t)$, let $\nu$ denote the unit
normal compatible with the orientation of $\Gamma_0$, and define
\begin{equation}
  H=-\operatorname{div}_\Gamma\nu,
  \qquad
  \Delta_\Gamma=\operatorname{div}_\Gamma\nabla_\Gamma.
\end{equation}
Let $V$ denote the normal velocity in the direction $\nu$. For a scalar field
$q$ on $\Gamma(t)$, let $\partial_t^\circ q$ denote its normal time
derivative. Equivalently, if $\bar q$ is a local space--time extension of
$q$, then
\begin{equation}
  \partial_t^\circ q
  =
  \left(
    \partial_t\bar q+V\nu\cdot\nabla\bar q
  \right)\big|_{\Gamma(t)}.
\end{equation}

Consider the system
\begin{equation}
  \boxed{
  \begin{aligned}
    V
    &=
    -\Delta_\Gamma
    \left[\left(1-\frac{q^2}{2}\right)H\right],\\[0.4em]
    \partial_t^\circ q
    &=
    \Delta_\Gamma q+qHV,\\[0.4em]
    (\Gamma,q)\big|_{t=0}
    &=
    (\Gamma_0,q_0).
  \end{aligned}}
  \label{eq:system}
\end{equation}

For $T>0$, define
\begin{align}
  \Erho(T)
  &:={W_p^1(0,T;L_p(\Sigma))}
    \cap L_p(0,T;W_p^4(\Sigma)),\\
  \Eq(T)
  &:={W_p^1(0,T;L_p(\Sigma))}
    \cap L_p(0,T;W_p^2(\Sigma)),
\end{align}
equipped with their natural sum norms.

Define the admissible phase space
\begin{equation}
  \mathcal O
  :=
  \left\{
    (\rho,u)\in
    \mathcal U\times W_p^{2-2/p}(\Sigma):
    \lVert u\rVert_{L^\infty(\Sigma)}<\sqrt2
  \right\}.
\end{equation}

\section{Problem Formalization}

Prove the following statements.

\begin{enumerate}
  \item There exists a maximal existence time
  \begin{equation}
    T_*\in(0,\infty]
  \end{equation}
  and a unique solution $(\Gamma(t),q(t))_{0\leq t<T_*}$ of
  \eqref{eq:system} such that
  \begin{equation}
    \Gamma(t)=X_{\rho(t)}(\Sigma),
    \qquad
    \widetilde q(t):=q(t)\circ X_{\rho(t)}.
  \end{equation}
  For every $T<T_*$, prove that
  \begin{equation}
    \rho\in\Erho(T),
    \qquad
    \widetilde q\in\Eq(T),
  \end{equation}
  and
  \begin{equation}
    \rho\in C\bigl([0,T];W_p^{4-4/p}(\Sigma)\bigr),
    \qquad
    \widetilde q\in C\bigl([0,T];W_p^{2-2/p}(\Sigma)\bigr),
  \end{equation}
  with
  \begin{equation}
    \rho(0)=\rho_0,
    \qquad
    \widetilde q(0)=\widetilde q_0.
  \end{equation}
  Moreover, prove that
  \begin{equation}
    \bigl(\rho(t),\widetilde q(t)\bigr)\in\mathcal O
    \qquad\text{for every }0\leq t<T_*.
  \end{equation}

  \item Prove continuous dependence on the initial data in
  \begin{equation}
    W_p^{4-4/p}(\Sigma)\times W_p^{2-2/p}(\Sigma)
  \end{equation}
  in the following precise sense. If
  \begin{equation}
    z_0=(\rho_0,\widetilde q_0)\in\mathcal O
  \end{equation}
  and $T<T_*(z_0)$, then there exists a neighborhood
  $\mathcal N\subset\mathcal O$ of $z_0$ such that
  \begin{equation}
    T_*(z_0')>T
    \qquad\text{for every }z_0'\in\mathcal N,
  \end{equation}
  and the solution map
  \begin{equation}
    z_0'
    \longmapsto
    \bigl(
      \rho(\,\cdot\,;z_0'),
      \widetilde q(\,\cdot\,;z_0')
    \bigr)
  \end{equation}
  is continuous from $\mathcal N$ into
  \begin{equation}
    \Erho(T)\times\Eq(T).
  \end{equation}

  Prove geometric uniqueness in the corresponding maximal-regularity class:
  any two solutions with the same initial oriented surface and the same
  initial scalar field determine the same evolving oriented surfaces and the
  same geometric scalar field.

  Consequently, prove uniqueness up to time-dependent reparametrization. More
  precisely, if $F_i(t,\cdot):M_i\to\mathbb R^3$, $i=1,2$, are two
  parametrized embedded solutions representing the same initial geometric
  data, and
  \begin{equation}
    u_i(t)=q_i(t)\circ F_i(t),
  \end{equation}
  then there exists a family of diffeomorphisms
  \begin{equation}
    \phi_t:M_2\longrightarrow M_1
  \end{equation}
  such that
  \begin{equation}
    F_2(t)=F_1(t)\circ\phi_t,
    \qquad
    u_2(t)=u_1(t)\circ\phi_t.
  \end{equation}

  \item The solution is instantly smooth. More precisely, for every
  \begin{equation}
    0<\tau<T'<T_*
  \end{equation}
  and every integer $m\geq0$, prove that
  \begin{equation}
    \rho
    \in
    W_p^1(\tau,T';W_p^m(\Sigma))
    \cap
    L_p(\tau,T';W_p^{m+4}(\Sigma))
  \end{equation}
  and
  \begin{equation}
    \widetilde q
    \in
    W_p^1(\tau,T';W_p^m(\Sigma))
    \cap
    L_p(\tau,T';W_p^{m+2}(\Sigma)).
  \end{equation}

  \item Assume that $T_*<\infty$ and that there exist
  \begin{equation}
    \rho_*\in W_p^{4-4/p}(\Sigma),
    \qquad
    \widetilde q_*\in W_p^{2-2/p}(\Sigma)
  \end{equation}
  such that
  \begin{equation}
    \rho(t)\longrightarrow\rho_*
    \quad\text{in }W_p^{4-4/p}(\Sigma),
  \end{equation}
  and
  \begin{equation}
    \widetilde q(t)\longrightarrow\widetilde q_*
    \quad\text{in }W_p^{2-2/p}(\Sigma)
  \end{equation}
  as $t\uparrow T_*$. Prove that if
  \begin{equation}
    \rho_*\in\mathcal U,
    \qquad
    \lVert\widetilde q_*\rVert_{L^\infty(\Sigma)}<\sqrt2,
  \end{equation}
  then there exists $\varepsilon>0$ such that the solution extends to
  \begin{equation}
    [0,T_*+\varepsilon).
  \end{equation}

  \item Deduce that the solution extends beyond every finite $T_*$ for which
  there exists $\delta>0$ such that
  \begin{equation}
    \inf_{0\leq t<T_*}
    \operatorname{dist}_{W_p^{4-4/p}(\Sigma)}
    \bigl(\rho(t),\partial\mathcal U\bigr)
    \geq\delta,
  \end{equation}
  \begin{equation}
    \sup_{0\leq t<T_*}
    \lVert\widetilde q(t)\rVert_{L^\infty(\Sigma)}
    \leq\sqrt2-\delta,
  \end{equation}
  and
  \begin{equation}
    \lVert\rho\rVert_{\Erho(T_*)}
    +
    \lVert\widetilde q\rVert_{\Eq(T_*)}
    <\infty.
  \end{equation}

  Equivalently, prove that if $T_*<\infty$ is the maximal existence time,
  then at least one of the following must occur:
  \begin{enumerate}
    \item
    \begin{equation}
      \operatorname{dist}_{W_p^{4-4/p}(\Sigma)}
      \bigl(\rho(t),\partial\mathcal U\bigr)
      \longrightarrow0
    \end{equation}
    along a sequence $t\uparrow T_*$;

    \item
    \begin{equation}
      \lVert\widetilde q(t)\rVert_{L^\infty(\Sigma)}
      \longrightarrow\sqrt2
    \end{equation}
    along a sequence $t\uparrow T_*$;

    \item
    \begin{equation}
      \lVert\rho\rVert_{\Erho(T_*)}
      +
      \lVert\widetilde q\rVert_{\Eq(T_*)}
      =\infty.
    \end{equation}
  \end{enumerate}

  \item Prove conservation of the total mass:
  \begin{equation}
    \int_{\Gamma(t)}q(t)\,dA
    =
    \int_{\Gamma_0}q_0\,dA
    \qquad\text{for every }0\leq t<T_*.
  \end{equation}

  Define the energy
  \begin{equation}
    E(t)
    :=
    \int_{\Gamma(t)}
    \left(1+\frac{q(t)^2}{2}\right)dA.
  \end{equation}
  Prove that
  \begin{equation}
    \frac{d}{dt}E(t)
    =
    -\int_{\Gamma(t)}
    \left(
      |\nabla_\Gamma q|^2
      +
      \left|
        \nabla_\Gamma
        \left[\left(1-\frac{q^2}{2}\right)H\right]
      \right|^2
    \right)dA
  \end{equation}
  for almost every $t\in(0,T_*)$.

  Equivalently, prove that for every $0\leq t<T_*$,
  \begin{align}
    E(t)
    &+
    \int_0^t\int_{\Gamma(s)}
    \Bigg(
      |\nabla_{\Gamma(s)}q|^2
      +
      \left|
        \nabla_{\Gamma(s)}
        \left[\left(1-\frac{q^2}{2}\right)H\right]
      \right|^2
    \Bigg)
    \,dA\,ds
    \notag\\
    &=E(0).
  \end{align}
\end{enumerate}


\end{document}
